\chapter{रैखिक अवकल समीकरण}\label{ch:b1:diffeq}

अवकल समीकरण पहले उच्चतर माध्यमिक खंड में मिले थे; यहाँ उनका उपचार पूरी
उपपत्तियों के साथ और अधिक व्यापकता में है: चर गुणांकों वाले प्रथम कोटि के
रैखिक समीकरण (जिन्हें अचरों के विचरण की विधि पूरी तरह हल कर देती है) और अचर
गुणांकों वाले द्वितीय कोटि के रैखिक समीकरण, जो दोलनों के आदर्श हैं। दोनों
स्थितियाँ एक ही संरचना दिखाती हैं: \emph{व्यापक हल $=$ एक विशिष्ट हल $+$
समांगी समीकरण का व्यापक हल}।

\section{प्रथम कोटि के रैखिक समीकरण}

\begin{definition}\label{def:b1:diffeq:linear1}
मान लीजिए $I$ एक अंतराल है और $a, b \colon I \to \R$ (या $\C$) संतत है।
समीकरण
\[
(E)\colon\quad y' + a(x)\,y = b(x),
\]
अज्ञात अवकलनीय फलन $y \colon I \to \R$ (या $\C$) में \emph{प्रथम कोटि का
रैखिक अवकल समीकरण}\index{अवकल समीकरण!प्रथम कोटि रैखिक} कहलाता है। समीकरण
$(H)\colon y' + a(x) y = 0$ उसका \emph{समांगी} समीकरण है।
\end{definition}

\begin{theorem}[समांगी समीकरण हल करना]\label{thm:b1:diffeq:homogeneous1}
मान लीजिए $A$ $I$ पर $a$ का प्रतिअवकलज है (वह विद्यमान है:
\cref{ch:b1:integration})। $I$ पर $(H)$ के हल ठीक ये फलन हैं
\[
y(x) = \lambda\, \eu^{-A(x)}, \qquad \lambda \in \R \text{ (या } \C).
\]
\end{theorem}

\begin{proof}
ये फलन हल हैं: $y' = -\lambda A' \eu^{-A} = -a y$। विलोमतः, मान लीजिए $y$
$(H)$ को हल करता है और $z(x) = y(x)\, \eu^{A(x)}$ रखिए। तब
\[
z' = y' \eu^{A} + y\, a\, \eu^{A} = (y' + ay)\, \eu^{A} = 0 ,
\]
अतः $z$ अंतराल $I$ पर अचर है, मान लीजिए $z = \lambda$: $y = \lambda
\eu^{-A}$। (तर्क पर ध्यान दीजिए: कोई हल नहीं छूटता, क्योंकि \emph{प्रत्येक}
हल घोषित रूप में लिखा जा चुका है।)
\end{proof}

\begin{example}[चर गुणांक वाला समांगी समीकरण]\label{ex:b1:diffeq:varcoeff}
$\R$ पर $y' + (\cos x)\,y = 0$ हल कीजिए। $a(x) = \cos x$ का एक प्रतिअवकलज
$A(x) = \sin x$ है, अतः हल हैं
\[
y(x) = \lambda\,\eu^{-\sin x}, \qquad \lambda \in \R .
\]
दो पाठ। प्रत्येक हल आवर्ती है (आवर्तकाल $2\pi$) और $\lambda = 0$ के अतिरिक्त
कभी लुप्त नहीं होता --- $\lambda$ का चिह्न सदा के लिए $y$ का चिह्न रहता है,
क्योंकि चरघातांकी शून्य को पार नहीं कर सकता। और $y(0) = y_0$ से होकर जाने
वाला हल $y_0\eu^{-\sin x}$ है: प्रत्येक आरंभिक बिंदु से होकर इस कुल का ठीक
एक वक्र, जो \cref{thm:b1:diffeq:voc}~(2) का एकविमीय चित्र है।
\end{example}

\begin{theorem}[अचरों का विचरण; कोशी समस्या]\label{thm:b1:diffeq:voc}
ऊपर के संकेतन के साथ:
\begin{enumerate}
  \item $I$ पर $(E)$ के हल ठीक ये हैं
        \[
        y(x) = \Bigl(\lambda + \int_{x_0}^{x} b(t)\, \eu^{A(t)} \dd t
        \Bigr)\, \eu^{-A(x)},
        \qquad \lambda \in \R,
        \]
        जहाँ $x_0 \in I$ नियत है। समतुल्य रूप से: $(H)$ का व्यापक हल और
        $(E)$ का एक विशिष्ट हल।
  \item प्रत्येक $x_0 \in I$ और $y_0$ के लिए \emph{कोशी समस्या}\index{कोशी
  समस्या} ``$(E)$ और $y(x_0) = y_0$'' का $I$ पर ठीक एक हल है।
\end{enumerate}
\end{theorem}

\begin{proof}
(1) \emph{अचरों का विचरण}\index{अचरों का विचरण} नामक विधि के अनुसार $y =
\mu(x)\, \eu^{-A(x)}$ रूप के हल खोजिए, जहाँ $\mu$ अवकलनीय है --- इससे
व्यापकता नहीं खोती, क्योंकि $I$ पर प्रत्येक फलन इस रूप में लिखा जा सकता है
($\mu = y\,\eu^{A}$)। प्रतिस्थापित करने पर,
\[
y' + ay = \mu' \eu^{-A} - \mu a \eu^{-A} + a\mu\eu^{-A}
= \mu'\, \eu^{-A},
\]
अतः $y$ $(E)$ को हल करता है यदि और केवल यदि $\mu'(x) = b(x)\,\eu^{A(x)}$,
अर्थात् यदि और केवल यदि किसी अचर $\lambda$ के लिए $\mu(x) = \lambda +
\int_{x_0}^x b(t)\eu^{A(t)}\dd t$ (अंतराल पर एक ही संतत फलन के दो प्रतिअवकलज
अचर से भिन्न होते हैं)।

(2) सूत्र में $y(x_0) = \lambda\,\eu^{-A(x_0)}$: शर्त $y(x_0) = y_0$
$\lambda = y_0 \eu^{A(x_0)}$ को अद्वितीय रूप से निर्धारित कर देती है।
\end{proof}

\begin{example}\label{ex:b1:diffeq:order1}
$I = \intoo{0}{+\infty}$ पर $y' + \dfrac{y}{x} = x^2$ हल कीजिए। यहाँ $a(x) =
\frac 1x$, $A(x) = \ln x$, $\eu^{-A(x)} = \frac 1x$। समांगी हल:
$\frac{\lambda}{x}$। \omterm{thm:b1:diffeq:voc}{अचरों का विचरण}: $\mu'(x) = x^2 \cdot x = x^3$, अतः $\mu
= \frac{x^4}{4} + \lambda$, और
\[
y(x) = \frac{x^3}{4} + \frac{\lambda}{x}, \qquad \lambda \in \R .
\]
आरंभिक शर्त $y(1) = 0$ के साथ: $\lambda = -\frac14$। \emph{जाँच:} $y' +
\frac yx = \frac{3x^2}{4} - \frac{\lambda}{x^2} + \frac{x^2}{4} +
\frac{\lambda}{x^2} = x^2$।
\end{example}

\begin{example}[अनुमान लगाना समाकलन से अच्छा है]\label{ex:b1:diffeq:guess}
$\R$ पर $y' + 2x\,y = x$ हल कीजिए। \omterm{thm:b1:diffeq:voc}{अचरों का विचरण} काम करता है ($A = x^2$,
$\mu' = x\,\eu^{x^2}$, $\mu = \frac12\eu^{x^2} + \lambda$), पर यह देख लेना
कि \emph{अचर} $y_p = \frac12$ समीकरण को हल करता है ($0 + 2x\cdot\frac12 =
x$), अधिक तेज़ है। समांगी हलों $\lambda\,\eu^{-x^2}$ के साथ:
\[
y(x) = \frac12 + \lambda\,\eu^{-x^2}, \qquad \lambda \in \R .
\]
जब $x \to \pm\infty$, तब प्रत्येक हल अत्यंत तेज़ी से $\frac12$ की ओर अभिसरित
होता है: अचर विशिष्ट हल एक \emph{साम्यावस्था} है जिससे सभी हल आ मिलते हैं।
सार: सामान्य विधि चलाने से पहले दस सेकंड कोई स्पष्ट विशिष्ट हल (अचर, एकपदी,
दाएँ पक्ष का गुणज) ढूँढ़ने में लगाइए; संरचना प्रमेय फिर काम पूरा कर देती है।
\end{example}

\begin{remark}[अंतराल महत्त्व रखते हैं]
यह प्रमेय ऐसे \emph{अंतराल} पर रहती है जहाँ $a$ और $b$ संतत हों। $\R^*$ पर
$y' + \frac yx = 0$ के लिए हल $\intoo{0}{+\infty}$ पर $\frac{\lambda}{x}$ और
$\intoo{-\infty}{0}$ पर $\frac{\mu}{x}$ हैं, जिनके अचर \emph{स्वतंत्र} हैं:
$0$ की विचित्रता को पार करके एक ही सूत्र के जुड़ने का कोई कारण नहीं है।
\end{remark}

\begin{example}[सम्मिश्र दायाँ पक्ष, दो वास्तविक उत्तर]\label{ex:b1:diffeq:complexrhs}
$y' - y = \cos x$ और $y' - y = \sin x$ को एक ही चाल में हल कीजिए। दाएँ पक्ष
$\eu^{\iu x}$ के साथ $\C$ में काम कीजिए: $y_p = c\,\eu^{\iu x}$ आज़माने पर
$c(\iu - 1)\eu^{\iu x} = \eu^{\iu x}$ मिलता है, अतः
\[
c = \frac1{\iu - 1} = \frac{-1 - \iu}2,
\qquad
y_p = -\frac{(1 + \iu)(\cos x + \iu\sin x)}2
= \frac{\sin x - \cos x}2 + \iu\,\frac{-\sin x - \cos x}2 .
\]
चूँकि समीकरण के गुणांक वास्तविक हैं, वास्तविक और काल्पनिक भाग अलग हो जाते
हैं: $\frac{\sin x - \cos x}2$ $y' - y = \cos x$ को हल करता है, और
$-\frac{\sin x + \cos x}2$ $y' - y = \sin x$ को (पहले की जाँच: अवकलज
$\frac{\cos x + \sin x}2$, में से फलन घटाने पर $\cos x$ मिलता है)। एक
सम्मिश्र पंक्ति ने अचरों के विचरण की दो दौड़ों की जगह ले ली --- वही
मितव्ययिता जिसे \cref{met:b1:diffeq:particular} द्वितीय कोटि के लिए
व्यवस्थित करता है, और \cref{ch:b1:complex} का बार-बार मिलने वाला लाभांश।
\end{example}

\section{अचर गुणांकों वाले द्वितीय कोटि के रैखिक समीकरण}

\begin{definition}\label{def:b1:diffeq:linear2}
मान लीजिए $a, b \in \R$ और $f \colon I \to \R$ संतत है। समीकरण
\[
(E)\colon\quad y'' + a\,y' + b\,y = f(x)
\]
\emph{अचर गुणांकों वाला द्वितीय कोटि का रैखिक समीकरण}\index{अवकल
समीकरण!द्वितीय कोटि रैखिक} कहलाता है; $(H)\colon y'' + ay' + by = 0$ उसका
समांगी समीकरण है, और $\chi(r) = r^2 + ar + b$ उसका \emph{अभिलक्षणिक
बहुपद}\index{अभिलक्षणिक बहुपद}।
\end{definition}

\begin{theorem}[समांगी हल]\label{thm:b1:diffeq:homogeneous2}
मान लीजिए $\Delta = a^2 - 4b$ $\chi$ का विविक्तकर है। $\R$ पर $(H)$ के
वास्तविक हल ये हैं:
\begin{enumerate}
  \item यदि $\Delta > 0$, तो दो वास्तविक मूलों $r_1 \neq r_2$ के साथ: $\;y =
  \lambda\, \eu^{r_1 x} + \mu\, \eu^{r_2 x}$;
  \item यदि $\Delta = 0$, तो द्विक मूल $r_0$ के साथ: $\;y = (\lambda + \mu
  x)\, \eu^{r_0 x}$;
  \item यदि $\Delta < 0$, तो मूलों $\alpha \pm \iu\omega$ ($\omega > 0$) के
  साथ: $\;y = \eu^{\alpha x} (\lambda \cos\omega x + \mu \sin\omega x)$;
\end{enumerate}
और हर स्थिति में $(\lambda, \mu)$ $\R^2$ पर चलते हैं।
\end{theorem}

\begin{proof}
पहले यह देखिए कि $r \in \C$ के लिए $x \mapsto \eu^{rx}$ $(H)$ को हल करता है
यदि और केवल यदि $\chi(r) = 0$ (प्रतिस्थापित कीजिए: $(r^2 + ar + b)\eu^{rx} =
0$)। इसीलिए चरघातांकी स्वाभाविक पहला अनुमान हैं: अवकलन $\eu^{rx}$ पर संख्या
$r$ से गुणन की भाँति काम करता है, अतः अवकल समीकरण संख्यात्मक समीकरण $\chi(r)
= 0$ बन जाता है --- पूरी विश्लेषणात्मक समस्या एक ही द्विघात के मूल खोजने में
सिमट जाती है।

मुख्य पद अज्ञात का ऐसा परिवर्तन है जो कोटि घटा देता है। मान लीजिए $r$ $\chi$
का (संभवतः सम्मिश्र) मूल है और $y = z\, \eu^{rx}$ लिखिए, जिससे व्यापकता नहीं
खोती। तब
\[
y'' + ay' + by
= \bigl(z'' + (2r + a) z' + \chi(r) z\bigr)\eu^{rx}
= \bigl(z'' + (2r + a) z'\bigr)\eu^{rx},
\]
अतः $(H)$ $u = z'$ के लिए \emph{प्रथम कोटि} का समीकरण $u' + (2r + a) u = 0$
बन जाता है।

\emph{स्थिति $\Delta \neq 0$:} $r = r_1$ चुनिए; तब $2r_1 + a = r_1 - r_2$
(क्योंकि $r_1 + r_2 = -a$)। \cref{thm:b1:diffeq:homogeneous1} से, किसी अचर
$c$ के लिए $z' = c\,\eu^{(r_2 - r_1)x}$; $\R$ पर समाकलन करने पर $\mu =
\frac{c}{r_2 - r_1}$ के साथ $z = \mu\, \eu^{(r_2 - r_1)x} + \lambda$, और
इसलिए $y = z\, \eu^{r_1 x} = \lambda\,\eu^{r_1x} + \mu\,\eu^{r_2x}$। जब
$\Delta < 0$, तब मूल $\alpha \pm \iu\omega$ हैं और सम्मिश्र हल $c_1, c_2 \in
\C$ के साथ $y = c_1\eu^{(\alpha+\iu\omega)x} + c_2\eu^{(\alpha-\iu\omega)x}$
हैं। इनमें से कौन वास्तविक-मान वाले हैं? चूँकि
$\conj{\eu^{(\alpha+\iu\omega)x}} = \eu^{(\alpha-\iu\omega)x}$, $y$ का
संयुग्मी $\conj{c_2}\, \eu^{(\alpha+\iu\omega)x} +
\conj{c_1}\,\eu^{(\alpha-\iu\omega)x}$ है, और सभी $x$ के लिए $y = \conj y$
$c_2 = \conj{c_1}$ पर बाध्य कर देता है (दोनों चरघातांकी रैखिकतः स्वतंत्र
हैं: दो बिंदुओं पर मान लीजिए, अथवा $\eu^{\alpha x}$ से भाग देकर $x=0$ पर
तुलना कीजिए)। वास्तविक $\lambda, \mu$ के साथ $c_1 = \frac{\lambda -
\iu\mu}2$ लिखने पर:
\[
y = 2\,\Re\Bigl(c_1\,\eu^{\alpha x}(\cos\omega x +
\iu\sin\omega x)\Bigr)
= \eu^{\alpha x}\bigl(\lambda\cos\omega x + \mu\sin\omega
x\bigr),
\]
और विलोमतः ऐसा हर फलन एक हल है (वास्तविक समीकरण के सम्मिश्र हल का वास्तविक
भाग): अतः वास्तविक हल-समष्टि वही है जो घोषित की गई थी।

\emph{स्थिति $\Delta = 0$:} $r = r_0$, $2r_0 + a = 0$, अतः $z'' = 0$: $z =
\lambda + \mu x$ और $y = (\lambda + \mu x)\eu^{r_0 x}$।
\end{proof}

\begin{example}[एक कोशी समस्या, आरंभ से अंत तक]\label{ex:b1:diffeq:cauchy2}
$y(0) = 0$, $y'(0) = 1$ के साथ $y'' - 3y' + 2y = 0$ हल कीजिए। \omterm{def:b1:diffeq:linear2}{अभिलक्षणिक
बहुपद} $r^2 - 3r + 2 = (r - 1)(r - 2)$ के वास्तविक मूल $1$ और $2$ हैं: व्यापक
हल $y = \lambda\eu^{x} + \mu\eu^{2x}$। दोनों शर्तें रैखिक निकाय देती हैं
\[
\lambda + \mu = 0, \qquad \lambda + 2\mu = 1 ,
\]
अतः $\mu = 1$, $\lambda = -1$:
\[
y(x) = \eu^{2x} - \eu^{x} .
\]
जाँच: $y(0) = 0$; $y' = 2\eu^{2x} - \eu^x$ के लिए $y'(0) = 1$; और $y'' - 3y'
+ 2y = (4 - 6 + 2)\eu^{2x} + (-1 + 3 - 2)\eu^x = 0$। उत्तर के आकार पर ध्यान
दीजिए: $-\infty$ के पास धीमी विधा $-\eu^x$ प्रभावी है; $+\infty$ के पास तेज़
विधा $\eu^{2x}$। हलों को भिन्न क्षय या वृद्धि दरों वाली विधाओं के अध्यारोपण
के रूप में पढ़ना लाभदायक आदत है --- सप्ताहांत समस्या का क्षणिक/स्थायी \omterm{thm:b1:logic:partition}{विभाजन}
इसी प्रकार सँभाला गया है।
\end{example}

\begin{omfigure}
\begin{tikzpicture}[xscale=1.35, yscale=1.55]
  \draw[->, thin] (0,0) -- (6.4,0) node[below] {$x$};
  \draw[->, thin] (0,-0.85) -- (0,1.25) node[left] {$y$};
  \draw[omDef, thick, domain=0:6.2, samples=400, smooth]
    plot (\x, {exp(-0.35*\x)*cos(4*\x r)});
  \draw[omThm, thick, domain=0:6.2, samples=120, smooth]
    plot (\x, {(1+1.8*\x)*exp(-1.1*\x)});
  \draw[omProp, thick, domain=0:6.2, samples=120, smooth]
    plot (\x, {1.15*exp(-0.45*\x) - 0.15*exp(-2.5*\x)});
  \node[omDef, font=\small] at (2.1,-0.62)
    {$\Delta < 0$: दोलन करता है};
  \node[omThm, font=\small] at (0.95,1.1) {$\Delta = 0$};
  \node[omProp, font=\small] at (4.9,0.38)
    {$\Delta > 0$: कोई दोलन नहीं};
\end{tikzpicture}

{\small क्षयशील हलों वाले $y'' + ay' + by = 0$ की तीन अवस्थाएँ: अवमंदित दोलन
(सम्मिश्र मूल), क्रांतिक वापसी (द्विक मूल), अति-अवमंदित क्षय (दो वास्तविक
मूल)। कौन-सी अवस्था होगी, यह केवल $\Delta = a^2 - 4b$ के चिह्न से पढ़ ली
जाती है --- कुछ भी हल करने से पहले।}
\end{omfigure}

\begin{theorem}[संरचना और कोशी समस्या]\label{thm:b1:diffeq:structure2}
\begin{enumerate}
  \item यदि $y_p$ $(E)$ का एक विशिष्ट हल है, तो $(E)$ के हल ठीक $y_p + y_h$
  हैं, जहाँ $y_h$ $(H)$ के हलों पर चलता है।
  \item (अध्यारोपण) यदि $y_1$ $y'' + ay' + by = f_1$ को हल करता है और $y_2$
  $y'' + ay' + by = f_2$ को, तो $y_1 + y_2$ दाएँ पक्ष $f_1 + f_2$ वाले
  समीकरण को हल करता है।
  \item सभी $x_0 \in I$ और $(y_0, y_0')$ के लिए \omterm{thm:b1:diffeq:voc}{कोशी समस्या} ``$(E)$, $y(x_0)
  = y_0$, $y'(x_0) = y_0'$'' का $I$ पर ठीक एक हल है। \emph{(अस्तित्व, एक
  विशिष्ट हल मान लेने पर; अद्वितीयता पूरी तरह।)}
\end{enumerate}
\end{theorem}

\begin{proof}
(1) $y \mapsto y'' + ay' + by$ की रैखिकता से $y$ $(E)$ को हल करता है यदि और
केवल यदि $y - y_p$ $(H)$ को हल करता है। (2) वही रैखिकता है।

(3) (1) से यह सिद्ध करना पर्याप्त है कि अचर $(\lambda, \mu)$ किसी भी आँकड़े
$(y_0, y_0')$ के अनुसार सदा, और अद्वितीय रूप से, समायोजित किए जा सकते हैं।
चर का स्थानांतरण करके $x_0 = 0$ मान लीजिए। \cref{thm:b1:diffeq:homogeneous2}
की स्थिति (1) में $y = \lambda\eu^{r_1x} + \mu\eu^{r_2x}$ देता है
\[
y(0) = \lambda + \mu, \qquad y'(0) = r_1\lambda + r_2\mu :
\]
जो $(\lambda, \mu)$ में ऐसा रैखिक निकाय है जिसका सारणिक $r_2 - r_1 \neq 0$
है; उसे स्पष्ट रूप से हल करने पर $\mu = \frac{y_0' - r_1y_0}{r_2 - r_1}$ और
$\lambda = y_0 - \mu$: ठीक एक हल। स्थिति (2) में $y(0) = \lambda$ और $y'(0)
= r_0\lambda + \mu$: निकाय त्रिभुजीय है जिसका सारणिक $1$ है, और $\lambda =
y_0$, $\mu = y_0' - r_0y_0$ से हल हो जाता है। स्थिति (3) में $y(0) =
\lambda$ और $y'(0) = \alpha\lambda + \omega\mu$: सारणिक $\omega \neq 0$, और
$\lambda = y_0$, $\mu = \frac{y_0' - \alpha y_0}\omega$ से हल। हर स्थिति में
\omterm{def:b1:logic:map}{प्रतिचित्रण} $(\lambda, \mu) \mapsto (y(x_0), y'(x_0))$ एक रैखिक एकैकी
आच्छादन है --- \cref{ch:b1:linmaps} की भाषा इस स्थिति-दर-स्थिति जाँच को एक
वाक्य में सिकोड़ देगी।
\end{proof}

\begin{method}[$f(x) = P(x)\,\eu^{\gamma x}$ के लिए विशिष्ट हल]\label{met:b1:diffeq:particular}
जब दायाँ पक्ष $P(x)\,\eu^{\gamma x}$ हो, जहाँ $P$ बहुपद है और $\gamma \in
\R$ (यह बहुपदों, चरघातांकियों और सम्मिश्र $\gamma$ या अध्यारोपण के द्वारा
$\cos$ तथा $\sin$ को समेट लेता है), तब इस रूप का विशिष्ट हल खोजिए
\[
y_p(x) = x^{m}\, Q(x)\, \eu^{\gamma x},
\qquad
m = \text{बहुपद } \chi \text{ के मूल के रूप में } \gamma \text{ की बहुलता}
\ (m = 0, 1 \text{ या } 2),
\]
जहाँ $Q$ $P$ के समान घात का बहुपद है, जिसके गुणांक प्रतिस्थापन और पहचान से
मिलते हैं। $f = K\cos\omega x$ (या $\sin$) के लिए दाएँ पक्ष $K\eu^{\iu\omega
x}$ के साथ हल कीजिए और वास्तविक (क्रमशः काल्पनिक) भाग लीजिए।
\end{method}

\begin{example}[अध्यारोपण काम पर]\label{ex:b1:diffeq:superposition}
$\R$ पर $y'' - y = \eu^{x} + 4$ हल कीजिए। समांगी: $\chi(r) = r^2 - 1$, मूल
$\pm1$, अतः $y_h = \lambda\eu^x + \mu\eu^{-x}$। दाएँ पक्ष को तोड़िए और हर
टुकड़े का उपचार विधि-बॉक्स से कीजिए। \emph{टुकड़ा $\eu^x$:} यहाँ $\gamma =
1$ $\chi$ का सरल मूल है, अतः $y_1 = c\,x\,\eu^x$ आज़माइए: तब $y_1'' - y_1 =
c(x + 2)\eu^x - cx\eu^x = 2c\,\eu^x$, जिससे $c = \frac12$ मिलता है।
\emph{टुकड़ा $4$:} $\gamma = 0$ मूल नहीं है; अचर $y_2 = -4$ काम करता है।
अध्यारोपण (\cref{thm:b1:diffeq:structure2}~(2)) से:
\[
y = \frac{x\,\eu^x}2 - 4 + \lambda\,\eu^x + \mu\,\eu^{-x},
\qquad (\lambda, \mu) \in \R^2 .
\]
ध्यान दीजिए कि दोनों टुकड़ों ने \emph{भिन्न} आकार माँगे ($m = 1$ बनाम $m =
0$): बहुलता की जाँच प्रत्येक घातांक पर अलग-अलग लगती है, और अनुमान लगाने से
पहले दाएँ पक्ष को तोड़ने का यही पूरा प्रयोजन है।
\end{example}

\begin{example}[बहुलता का नियम काम पर]\label{ex:b1:diffeq:multiplicity}
$\R$ पर $y'' + y' = x$ हल कीजिए। दायाँ पक्ष $P(x) = x$ के साथ $P(x)\eu^{0
\cdot x}$ है, और $\gamma = 0$ $\chi(r) = r^2 + r = r(r + 1)$ का \emph{सरल}
मूल है: अतः $m = 1$, और सही अनुमान $y_p = x\,(\alpha x + \beta) = \alpha x^2
+ \beta x$ है, जो $P$ से एक घात ऊपर है। प्रतिस्थापित करने पर:
\[
y_p'' + y_p' = 2\alpha + (2\alpha x + \beta)
= 2\alpha x + (2\alpha + \beta) ,
\]
और $x$ से पहचान करने पर $\alpha = \frac12$, $\beta = -1$ मिलते हैं: $y_p =
\frac{x^2}2 - x$। व्यापक हल: $y = \frac{x^2}2 - x + \lambda +
\mu\,\eu^{-x}$। यदि हमने $y_p = \alpha x + \beta$ का अनुमान लगाया होता
(बहुलता की उपेक्षा करके), तो प्रतिस्थापन $y_p'' + y_p' = \alpha$ देता, जो
अचर है --- $\alpha, \beta$ का कोई चयन $x$ से मेल नहीं खा सकता, और यह विफलता
संरचनात्मक है: अचर तो पहले से ही समांगी समीकरण को हल करते हैं, अतः वे बाएँ
पक्ष के लिए अदृश्य हैं। गुणनखंड $x^m$ ठीक इसीलिए है कि समांगी हल-समष्टि से
बाहर निकला जा सके।
\end{example}

\begin{remark}[तीस सेकंड की बीमा-पॉलिसी]\label{rem:b1:diffeq:check}
इस अध्याय का हर हल किया गया समीकरण प्रतिस्थापन की जाँच पर समाप्त होता है, और
यह सजावट नहीं है। अवकल समीकरण की संगणना अनेक छोटे पदों को जोड़ती है (एक
प्रतिअवकलज, एक गुणनफल नियम, दो अचर), और एक ही चिह्न की भूल अदृश्य रूप से आगे
बढ़ती जाती है; अंतिम सूत्र को समीकरण में वापस रखने पर एक अवकलन के मूल्य पर
लगभग सारी भूलें पकड़ में आ जाती हैं। यह प्रतिवर्त तीन परतों में विकसित
कीजिए: अकेले \emph{विशिष्ट हल} की जाँच कीजिए (समांगी भाग तो वैसे भी कट जाता
है), पूरे हल पर \emph{आरंभिक शर्तों} की जाँच कीजिए, और जब कोई प्राचल हो तब
किसी \emph{अपकर्षी मान} की जाँच कीजिए (क्या व्यापक $\Omega$ का सूत्र $\Omega
= 0$ पर ज्ञात उत्तर दुबारा दे देता है?)। इस आदत की कीमत आधा मिनट है; और वह
``शायद ठीक'' को ``सत्यापित'' में बदल देती है।
\end{remark}

\begin{remark}[सामान्य भूलें]\label{rem:b1:diffeq:pitfalls}
\begin{enumerate}
  \item \emph{पहले मानकीकरण कीजिए।} सूत्र मानकर चलते हैं कि समीकरण $y' +
  a(x)y = b(x)$ है --- $y'$ पर गुणांक $1$ के साथ। $xy' - 2y = x^3$ के लिए
  $a$ और $b$ की पहचान करने से पहले $x$ से भाग दीजिए ($0$ से बचने वाले अंतराल
  पर), जैसा \cref{exo:b1:diffeq:2} में है।
  \item \emph{प्रति विमा एक अचर, और वह अंत में नियत।} प्रथम कोटि का व्यापक
  हल एक अचर रखता है, द्वितीय कोटि का दो; आरंभिक शर्तें \emph{पूरे} हल $y_p +
  y_h$ पर लगाई जाती हैं, कभी अकेले $y_h$ पर नहीं --- $y_p$ जोड़ने से पहले
  उन्हें लगा देना सबसे बार-बार होने वाली संरचनात्मक भूल है।
  \item \emph{बहुलता का ध्यान रखिए।} जो विशिष्ट-हल अनुमान समांगी समीकरण को
  हल करता हो, वह बाएँ पक्ष के लिए अदृश्य है; \cref{met:b1:diffeq:particular}
  का गुणनखंड $x^m$ वैकल्पिक नहीं है (\cref{ex:b1:diffeq:multiplicity})।
  \item \emph{अंतराल उत्तर का हिस्सा हैं।} हल उन्हीं अंतरालों पर रहते हैं
  जहाँ गुणांक संतत हों; किसी विचित्रता को पार करके जोड़ने से मिथ्या अचर बन
  सकते हैं (\cref{exo:b1:diffeq:12}) या अद्वितीयता नष्ट हो सकती है। ``$\R^*$
  पर हल कीजिए'' का अर्थ है दो स्वतंत्र समस्याएँ।
\end{enumerate}
\end{remark}

\begin{example}[अनुनाद से हटकर प्रणोदन]\label{ex:b1:diffeq:offresonance}
$\R$ पर $y'' + 4y = \sin x$ हल कीजिए। स्वाभाविक आवृत्ति $2$, प्रणोदन आवृत्ति
$1$: चूँकि $\iu$ $\chi(r) = r^2 + 4$ का मूल \emph{नहीं} है, बहुलता $m = 0$
है और एक सादा ज्यावक्रीय फलन पर्याप्त है। $y_p = \alpha\sin x$ आज़माइए
(कोज्या की आवश्यकता नहीं: समीकरण में $y'$ वाला पद नहीं है, और $\sin$ $\sin$
को फिर से उत्पन्न कर देता है):
\[
y_p'' + 4y_p = -\alpha\sin x + 4\alpha\sin x = 3\alpha\sin x ,
\]
अतः $\alpha = \frac13$ और व्यापक हल है
\[
y = \frac{\sin x}3 + \lambda\cos 2x + \mu\sin 2x .
\]
प्रत्येक हल परिबद्ध रहता है: आवृत्तियों $1$ (प्रणोदित) और $2$ (स्वाभाविक) के
दो दोलनों का अध्यारोपण। अगले उदाहरण से तुलना कीजिए, जहाँ स्वाभाविक आवृत्ति
\emph{पर} प्रणोदन उत्तर का आकार ही बदल देता है।
\end{example}

\begin{example}[एक प्रणोदित दोलन]\label{ex:b1:diffeq:oscillation}
$y'' + y = \cos x$, $y(0) = 0$, $y'(0) = 0$ हल कीजिए।

\emph{समांगी:} $\chi(r) = r^2 + 1$, मूल $\pm\iu$: $y_h = \lambda\cos x + \mu
\sin x$।

\emph{विशिष्ट:} दायाँ पक्ष $\Re(\eu^{\iu x})$ है, जहाँ $\gamma = \iu$ $\chi$
का सरल मूल है: $z_p = c\, x\, \eu^{\iu x}$ आज़माइए ($c \in \C$)। तब $z_p'' +
z_p = c\,(2\iu)\eu^{\iu x}$, जो $c = \frac{1}{2\iu} = -\frac\iu2$ के लिए
$\eu^{\iu x}$ के बराबर है। अतः $z_p = -\frac{\iu}{2} x (\cos x + \iu \sin
x)$ और $y_p = \Re(z_p) = \frac{x \sin x}{2}$।

\emph{व्यापक हल:} $y = \frac{x\sin x}{2} + \lambda\cos x + \mu\sin x$।
शर्तें: $y(0) = \lambda = 0$; $y' = \frac{\sin x + x\cos x}{2} + \mu\cos x$,
अतः $y'(0) = \mu = 0$। उत्तर: $y = \frac{x\sin x}{2}$ --- ऐसा दोलन जिसका
आयाम रैखिक रूप से बढ़ता है: यही \emph{अनुनाद}\index{अनुनाद} की परिघटना है,
जो निकाय को उसकी स्वाभाविक आवृत्ति पर प्रणोदित करने से होती है।
\end{example}

\begin{omfigure}
\begin{tikzpicture}
\begin{axis}[omaxis, width=11cm, height=5.2cm,
    xmin=0, xmax=25, ymin=-13, ymax=13,
    xtick={0,5,10,15,20,25}, ytick=\empty]
  \addplot[omThm, thick, domain=0:25, samples=400] {x*sin(deg(x))/2};
  \addplot[dashed, gray, domain=0:25] {x/2};
  \addplot[dashed, gray, domain=0:25] {-x/2};
\end{axis}
\end{tikzpicture}

{\small \omterm{ex:b1:diffeq:oscillation}{अनुनाद}: $y'' + y = \cos x$ का हल $y = \frac{x \sin x}{2}$ रेखाओं $y
= \pm\frac x2$ (बिंदुकित) के बीच, निरंतर बढ़ते आयाम के साथ दोलन करता है।}
\end{omfigure}

\begin{remark}[मध्यांतर: रैखिकता एक ज्यामिति है]\label{rem:b1:diffeq:interlude}
इस अध्याय के हर \omterm{def:b1:logic:sets}{हल-समुच्चय} के आकार पर फिर दृष्टि डालिए: एक विशेष हल, और उसके
साथ समांगी हलों की ऐसी समष्टि जिसमें एक मुक्त अचर (प्रथम कोटि) या दो
(द्वितीय कोटि) हैं। रैखिक बीजगणित के अध्याय
(\cref{ch:b1:vspaces,ch:b1:findim,ch:b1:linmaps}) ठीक-ठीक शब्दावली देंगे:
\omterm{def:b1:logic:map}{प्रतिचित्रण} $L(y) = y'' + ay' + by$ \emph{रैखिक} है, उसके समांगी हल $L$ की
\emph{अष्टि} बनाते हैं, जो ऐसी सदिश समष्टि है जिसकी \emph{विमा} समीकरण की
कोटि के बराबर है --- ``प्रति कोटि एक अचर'' की ईमानदार सामग्री यही है --- और
$L(y) = f$ का \omterm{def:b1:logic:sets}{हल-समुच्चय} एक \emph{ऐफ़ीन उपसमष्टि} है, अर्थात् अष्टि का
स्थानांतरण। यहाँ तक कि \cref{thm:b1:diffeq:structure2} का कोशी \omterm{def:b1:logic:map}{प्रतिचित्रण}
$(\lambda, \mu) \mapsto (y(x_0), y'(x_0))$ भी दो समतलों के बीच रैखिक एकैकी
आच्छादन है, अर्थात् एक व्युत्क्रमणीय $2 \times 2$ निकाय
(\cref{ch:b1:matrices})। इस अध्याय में कुछ भी दोबारा करने की आवश्यकता नहीं
होगी --- केवल उसका नाम बदलना होगा, और यह नामकरण रैखिक बीजगणित के लिए
सर्वोत्तम संभव तैयारी है: वहाँ की हर अमूर्त परिभाषा यहाँ पहले ही अपनी रोटी
कमा चुकी है।
\end{remark}

\begin{remark}[यह अध्याय कहाँ काम आता है]\label{rem:b1:diffeq:whereused}
संरचना प्रमेय --- $(E)$ के हल ``एक विशिष्ट हल और $(H)$ के हल'' बनाते हैं ---
उस प्रतिरूप की पहली उपस्थिति है जिसे \cref{ch:b1:vspaces,ch:b1:linmaps} नाम
देगा: $(H)$ का \omterm{def:b1:logic:sets}{हल-समुच्चय} रैखिक \omterm{def:b1:logic:map}{प्रतिचित्रण} $y \mapsto y'' + ay' + by$ की
\emph{अष्टि} है, और $(E)$ का \omterm{def:b1:logic:sets}{हल-समुच्चय} उसका ऐफ़ीन स्थानांतरण। \omterm{def:b1:diffeq:linear2}{अभिलक्षणिक
बहुपद} \cref{ch:b1:matrices} में आव्यूह के \omterm{def:b1:diffeq:linear2}{अभिलक्षणिक बहुपद} के रूप में फिर
प्रकट होता है: द्वितीय कोटि का समीकरण वेश बदला हुआ $2 \times 2$ प्रथम कोटि
का निकाय है, और स्नातक वर्ष 2 का खंड इस दृष्टिकोण को व्यवस्थित करता है।
अचरों के विचरण की माँगी हुई समाकलों की आपूर्ति \cref{ch:b1:integration} करता
है, और नीचे दी गई सप्ताहांत समस्या --- प्रणोदित अवमंदित दोलक --- विज्ञानों
के हर दोलन-प्रश्न का आदर्श स्थिति है, परिपथों से लेकर झूलते पुलों तक।
\end{remark}

\section{अभ्यास}

\begin{exercise}[$\star$]\label{exo:b1:diffeq:1}
$\R$ पर हल कीजिए: $\;y' + 2y = \eu^{3x}$; फिर \omterm{thm:b1:diffeq:voc}{कोशी समस्या} $y(0) = 1$।
\end{exercise}

\begin{exercise}[$\star$]\label{exo:b1:diffeq:2}
$\intoo{0}{+\infty}$ पर हल कीजिए: $\;x y' - 2y = x^3$ \emph{(पहले समीकरण को
मानकीकृत रूप में रखिए)}।
\end{exercise}

\begin{exercise}[$\star$]\label{exo:b1:diffeq:3}
$\R$ पर वास्तविक व्यापक हल देते हुए हल कीजिए: $\;y'' - 3y' + 2y = 0$; $\;y''
+ 4y' + 4y = 0$; $\;y'' - 2y' + 5y = 0$।
\end{exercise}

\begin{exercise}[$\star$]\label{exo:b1:diffeq:4}
$\R$ पर $y'' - y = x^2$ हल कीजिए, फिर \omterm{thm:b1:diffeq:voc}{कोशी समस्या} $y(0) = 0$, $y'(0) = 1$।
\end{exercise}

\begin{exercise}[$\star\star$]\label{exo:b1:diffeq:5}
$\intoo{-\frac\pi2}{\frac\pi2}$ पर हल कीजिए: $\;y' + y\tan x = \sin 2x$।
\end{exercise}

\begin{exercise}[$\star\star$]\label{exo:b1:diffeq:6}
$\R$ पर $y'' - 4y' + 3y = (2x + 1)\,\eu^{x}$ हल कीजिए। \emph{(बहुलता का
ध्यान रखिए: क्या $1$ \omterm{def:b1:diffeq:linear2}{अभिलक्षणिक बहुपद} का मूल है?)}
\end{exercise}

\begin{exercise}[$\star\star$]\label{exo:b1:diffeq:7}
$\R$ पर $y'' + 4y = \sin 2x + x$ हल कीजिए \emph{(अध्यारोपण; प्रत्येक दाएँ
पक्ष का उपचार अलग-अलग कीजिए)}।
\end{exercise}

\begin{exercise}[$\star\star$]\label{exo:b1:diffeq:8}
$T_0 = 80\,^\circ$C तापमान वाली कॉफ़ी का प्याला $20\,^\circ$C के कमरे में
रखा है। न्यूटन का शीतलन नियम $k > 0$ के साथ $T' = -k\,(T - 20)$ कहता है।
$T(t)$ के लिए हल कीजिए, और यह देखते हुए कि $10$ मिनट बाद कॉफ़ी $50\,^\circ$C
पर है, बताइए कि वह $25\,^\circ$C पर कब पहुँचेगी।
\end{exercise}

\begin{exercise}[$\star\star\star$]\label{exo:b1:diffeq:9}
(अवमंदित दोलक) $\varepsilon \geq 0$ के लिए $y'' + 2\varepsilon y' + y = 0$
पर विचार कीजिए।
\begin{enumerate}
  \item $\varepsilon \in \intco{0}{1}$, $\varepsilon = 1$ और $\varepsilon >
  1$ के लिए हल कीजिए।
  \item दिखाइए कि $\varepsilon > 0$ के लिए प्रत्येक हल $+\infty$ पर $0$ की
  ओर जाता है, और $\varepsilon = 0$ के लिए अशून्य हल ऐसा नहीं करते।
  \item $\varepsilon \in \intoo{0}{1}$ के लिए दिखाइए कि किसी अशून्य हल के
  शून्य नियमित अंतराल पर हैं, जिनका अंतराल $\frac{\pi}{\sqrt{1 -
  \varepsilon^2}}$ है।
\end{enumerate}
\end{exercise}

\begin{exercise}[$\star\star\star$]\label{exo:b1:diffeq:10}
ऐसे सभी दो बार अवकलनीय फलन $f \colon \R \to \R$ खोजिए कि
\[
\forall x, y \in \R, \qquad f(x + y) + f(x - y) = 2 f(x) f(y),
\]
जहाँ $f(0) \ne 0$ और $f$ अचर नहीं है। \emph{संकेत: $y$ नियत कीजिए, $0$ पर
$x$ के सापेक्ष दो बार अवकलन कीजिए; दिखाइए कि किसी अचर $c$ के लिए $f(0) = 1$
और $f'' = c f$; फिर $c$ के चिह्न के अनुसार हल कीजिए और जाँचिए कि कौन-से हल
फलनीय समीकरण को संतुष्ट करते हैं।}
\end{exercise}

\begin{exercise}[$\star\star$]\label{exo:b1:diffeq:11}
(ऑयलर समीकरण) $\intoo{0}{+\infty}$ पर $x^2 y'' - x y' + y = 0$ हल कीजिए।
\emph{संकेत: $z(t) = y(\eu^t)$ रखिए, अर्थात् $x = \eu^t$ प्रतिस्थापित कीजिए,
और दिखाइए कि $z$ अचर गुणांकों वाले रैखिक समीकरण को संतुष्ट करता है।}
\end{exercise}

\begin{exercise}[$\star\star\star$]\label{exo:b1:diffeq:12}
अज्ञात अवकलनीय फलन $y \colon \R \to \R$ के साथ पूरी वास्तविक रेखा पर समीकरण
$x\,y' = 2y$ पर विचार कीजिए।
\begin{enumerate}
  \item $\intoo{0}{+\infty}$ पर और $\intoo{-\infty}{0}$ पर हल कीजिए।
  \item दिखाइए कि \emph{किन्हीं भी} अचरों $a, b \in \R$ के लिए वह फलन जो $x
  \geq 0$ के लिए $ax^2$ के बराबर हो और $x < 0$ के लिए $bx^2$ के बराबर, $\R$
  पर अवकलनीय है और सर्वत्र समीकरण को हल करता है।
  \item निष्कर्ष निकालिए कि $\R$ पर \omterm{def:b1:logic:sets}{हल-समुच्चय} दो-प्राचल वाला कुल है, और
  समझाइए कि यह \cref{thm:b1:diffeq:voc} की अद्वितीयता का विरोध क्यों नहीं
  करता।
\end{enumerate}
\end{exercise}

\section{समस्या: प्रणोदित अवमंदित दोलक}

\begin{problem}\label{pb:b1:diffeq:1}
एक ही समीकरण श्यान माध्यम में स्प्रिंग पर लगे द्रव्यमान, RLC परिपथ के आवेश,
और हवा में झूलती इमारत --- तीनों को नियंत्रित करता है:
\[
(E_\Omega)\colon\quad
x'' + 2\lambda x' + \omega_0^2\,x = A\cos(\Omega t),
\]
जहाँ $\lambda \geq 0$ अवमंदन है, $\omega_0 > 0$ स्वाभाविक आवृत्ति है, और $A
> 0$, $\Omega > 0$ प्रणोदन के आयाम तथा आवृत्ति हैं। यह समस्या उसका पूरा
व्यवहार निकाल लेती है: क्षणिक भागों का क्षय, अद्वितीय आवर्ती स्थायी अवस्था,
\omterm{ex:b1:diffeq:oscillation}{अनुनाद} वक्र और उसकी तीक्ष्णता (\emph{\omterm{pb:b1:diffeq:1}{गुणता कारक}}), अनवमंदित स्थिति के
\omterm{pb:b1:diffeq:1}{विस्पंद}, और वह ऊर्जा-संतुलन जो दोलन को बनाए रखता
है।\index{अनुनाद}\index{दोलक} जब तक अन्यथा न कहा जाए, $0 < \lambda <
\omega_0$ (अल्प-अवमंदित अवस्था) और हम $\omega_d = \sqrt{\omega_0^2 -
\lambda^2}$ लिखते हैं।

\medskip \noindent\textbf{भाग I --- मुक्त दोलक।} यहाँ $A = 0$।
\begin{enumerate}
  \item $0 < \lambda < \omega_0$ के लिए और $\lambda = 0$ के लिए समांगी
  समीकरण $(H)$ हल कीजिए। (अवस्थाएँ $\lambda \geq \omega_0$
  \cref{exo:b1:diffeq:9} में सँभाली गई थीं; उन्हें उद्धृत कीजिए।)
  \item दिखाइए कि प्रत्येक $\lambda > 0$ के लिए $(H)$ के सभी हल $+\infty$ पर
  $0$ की ओर जाते हैं --- तीनों अवस्थाओं में।
  \item $(H)$ के किसी हल पर ऊर्जा $\mathcal E(t) = \frac12 x'(t)^2 +
  \frac12\omega_0^2\,x(t)^2$ परिभाषित कीजिए। दिखाइए $\mathcal E'(t) =
  -2\lambda\,x'(t)^2 \leq 0$, और उससे (बिना कुछ हल किए) निकालिए कि प्रत्येक
  $\lambda \geq 0$ के लिए \omterm{thm:b1:diffeq:voc}{कोशी समस्या} ``$(H)$, $x(t_0) = x'(t_0) = 0$'' का
  केवल शून्य हल है।
  \item $0 < \lambda < \omega_0$ के लिए अशून्य हल को $x(t) =
  R\,\eu^{-\lambda t}\cos(\omega_d t - \varphi)$ लिखिए और मान लीजिए $T_d =
  \frac{2\pi}{\omega_d}$ छद्म-आवर्तकाल है। दिखाइए कि $x(t + T_d) =
  \eu^{-\lambda T_d}\,x(t)$: प्रत्येक झूला पिछले झूले का अचर गुणक
  $\eu^{-\delta}$ से सिकुड़ा हुआ रूप है, $\delta =
  \frac{2\pi\lambda}{\omega_d}$ (\emph{लघुगणकीय ह्रासांक})। $\omega_0 = 1$,
  $\lambda = 0.1$ के लिए $\delta$ संगणित कीजिए।
  \item \emph{\omterm{pb:b1:diffeq:1}{गुणता कारक}} $Q = \dfrac{\omega_0}{2\lambda}$\index{गुणता कारक}
  परिभाषित कीजिए। दिखाइए कि समय $\frac1\lambda$ (एक आयाम-$\eu$-गुणन) के बाद
  दोलक $\frac{\omega_d}{2\pi\lambda}$ छद्म-आवर्तकाल पूरे कर चुका होता है, जो
  दुर्बल अवमंदन ($\lambda \ll \omega_0$) के लिए लगभग $\frac Q\pi$ है: \omterm{pb:b1:diffeq:1}{गुणता
  कारक} $\pi$ तक यही गिनता है कि आयाम के $\eu$ से घटने से पहले कितने दोलन बच
  रहे।
\end{enumerate}

\medskip \noindent\textbf{भाग II --- स्थायी अवस्था।} अब $A > 0$ और $\lambda
> 0$।
\begin{enumerate}[resume]
  \item $z \in \C$ (\cref{met:b1:diffeq:particular}) के साथ
  $z\,\eu^{\iu\Omega t}$ के वास्तविक भाग के रूप में विशिष्ट हल खोजिए। दिखाइए
  कि यह इनके साथ काम करता है
        \[
        z = \frac{A}{\omega_0^2 - \Omega^2 + 2\iu\lambda\Omega} .
        \]
  \item उससे स्थायी अवस्था को आयाम--कला रूप में निकालिए: $x_p(t) =
  R(\Omega)\cos\bigl(\Omega t - \varphi(\Omega)\bigr)$, जहाँ
        \[
        R(\Omega) = \frac{A}{\sqrt{(\omega_0^2 - \Omega^2)^2 +
        4\lambda^2\Omega^2}},
        \qquad
        \tan\varphi = \frac{2\lambda\Omega}{\omega_0^2 -
        \Omega^2},
        \quad \varphi \in \intoo0\pi .
        \]
  \item दोनों चरम अवस्थाओं की व्याख्या कीजिए: $\Omega \to 0^+$ पर $R$ और
  $\varphi$ की सीमाएँ संगणित कीजिए (अर्ध-स्थैतिक अनुक्रिया $A/\omega_0^2$,
  कला $0$) और $\Omega \to +\infty$ पर भी ($R \sim A/\Omega^2 \to 0$, कला
  $\to \pi$: द्रव्यमान अत्यधिक तेज़ प्रणोदन के विपरीत चलता है)।
  \item दिखाइए कि $(E_\Omega)$ का \emph{प्रत्येक} हल $x_p$ और $(H)$ के किसी
  हल का योग है, अतः आरंभिक शर्तें जो भी हों, वह $t \to +\infty$ पर स्थायी
  अवस्था $x_p$ की ओर अभिसरित होता है: क्षणिक भाग के मर जाने पर दोलक को यह
  स्मरण नहीं रहता कि उसका आरंभ कैसे हुआ था।
  \item दिखाइए कि $(E_\Omega)$ का \emph{एकमात्र} आवर्ती हल $x_p$ है।
  \item एक \omterm{thm:b1:diffeq:voc}{कोशी समस्या} अंत तक कीजिए: $x(0) = x'(0) = 0$ के साथ $x'' + 2x' +
  2x = \cos t$ के लिए दिखाइए कि हल है
        \[
        x(t) = \frac{\cos t + 2\sin t}5
        - \eu^{-t}\,\frac{\cos t + 3\sin t}5 ,
        \]
        और क्षणिक तथा स्थायी भागों की पहचान कीजिए।
\end{enumerate}

\medskip \noindent\textbf{भाग III --- \omterm{ex:b1:diffeq:oscillation}{अनुनाद} वक्र।} $\intoo0{+\infty}$ पर
$\Omega \mapsto R(\Omega)$ का अध्ययन।
\begin{enumerate}[resume]
  \item $u = \Omega^2$ और $g(u) = (\omega_0^2 - u)^2 + 4\lambda^2 u$ रखकर
  दिखाइए: यदि $2\lambda^2 < \omega_0^2$, तो $R$ \emph{\omterm{ex:b1:diffeq:oscillation}{अनुनाद} आवृत्ति}
  $\Omega_r = \sqrt{\omega_0^2 - 2\lambda^2}$ पर कठोर अधिकतम प्राप्त करता
  है, जहाँ
        \[
        R_{\max} = R(\Omega_r)
        = \frac{A}{2\lambda\sqrt{\omega_0^2 - \lambda^2}} .
        \]
  \item दिखाइए कि $\dfrac{R_{\max}}{R(0)} = Q\,\bigl(1 -
  \tfrac{\lambda^2}{\omega_0^2}\bigr)^{-1/2} \geq Q$: \omterm{ex:b1:diffeq:oscillation}{अनुनाद} पर प्रणोदन
  (मूलतः) \omterm{pb:b1:diffeq:1}{गुणता कारक} से प्रवर्धित होता है।
  \item सिद्ध कीजिए कि \emph{वेग} का आयाम $V(\Omega) = \Omega\,R(\Omega)$
  ठीक $\Omega = \omega_0$ पर अधिकतम है ($\Omega_r$ पर नहीं), और यह कि वहाँ
  कला $\varphi(\omega_0) = \frac\pi2$ है: $\Omega = \omega_0$ पर वेग बल के
  ठीक समान कला में है।
  \item (पट्ट-चौड़ाई) $g(u) = 2\,g(u_r)$ यथार्थ रूप से हल कीजिए, जहाँ $u_r =
  \omega_0^2 - 2\lambda^2$, और उससे निकालिए कि $R = R_{\max}/\sqrt2$ वाली
  दोनों आवृत्तियाँ $\Omega_\pm$ $\Omega_+^2 - \Omega_-^2 =
  4\lambda\sqrt{\omega_0^2 - \lambda^2}$ को संतुष्ट करती हैं; निष्कर्ष
  निकालिए कि दुर्बल अवमंदन के लिए पट्ट-चौड़ाई $\Omega_+ - \Omega_- \approx
  2\lambda$ है, अर्थात् $Q \approx \frac{\omega_0}{\Omega_+ - \Omega_-}$:
  तीक्ष्ण अनुनाद-शिखर उच्च-$Q$ निकायों के होते हैं।
  \item $\omega_0 = 1$, $\lambda = 0.05$ ($Q = 10$), $A = 1$ के लिए
  संख्यात्मक चित्र: $\Omega_r$, $R_{\max}$, स्थैतिक अनुक्रिया $R(0)$ और
  सन्निकट पट्ट-चौड़ाई संगणित कीजिए।
  \item दिखाइए कि यदि $2\lambda^2 \geq \omega_0^2$, तो $R$
  $\intoo0{+\infty}$ पर कठोरतः ह्रासमान है: भारी अवमंदन वाले निकायों में
  अनुनाद-शिखर होता ही नहीं।
\end{enumerate}

\medskip \noindent\textbf{भाग IV --- अवमंदन के बिना: \omterm{pb:b1:diffeq:1}{विस्पंद} और \omterm{ex:b1:diffeq:oscillation}{अनुनाद}।}
यहाँ $\lambda = 0$।
\begin{enumerate}[resume]
  \item $\Omega \neq \omega_0$ के लिए $x'' + \omega_0^2 x = A\cos(\Omega t)$
  का व्यापक हल खोजिए।
  \item \omterm{thm:b1:diffeq:voc}{कोशी समस्या} $x(0) = x'(0) = 0$ हल कीजिए और उत्तर को गुणनफल-रूप में
  बदलिए
        \[
        x(t) = \frac{2A}{\omega_0^2 - \Omega^2}\,
        \sin\Bigl(\frac{(\omega_0 - \Omega)t}2\Bigr)
        \sin\Bigl(\frac{(\omega_0 + \Omega)t}2\Bigr) .
        \]
  \item $\omega_0$ के निकट $\Omega$ के लिए इस गुणनफल को आवृत्ति
  $\frac{\omega_0 + \Omega}2$ के तीव्र दोलन के रूप में पढ़िए, जो आवृत्ति
  $\frac{\abs{\omega_0 - \Omega}}2$ के मंद अन्वालोप से मॉडुलित है: यही
  \emph{विस्पंद}\index{विस्पंद} हैं। अन्वालोप का आवर्तकाल और अधिकतम आयाम
  दीजिए, और ध्यान दीजिए कि $\Omega \to \omega_0$ होने पर दोनों कैसे फट पड़ते
  हैं।
  \item $t$ नियत कीजिए और प्रश्न 19 के सूत्र में $\Omega \to \omega_0$
  लीजिए: दिखाइए कि सीमा है
        \[
        x_\infty(t) = \frac{A\,t\,\sin(\omega_0 t)}{2\omega_0} ,
        \]
        और सीधे जाँचिए कि $x_\infty$ $x(0) = x'(0) = 0$ के साथ अनुनादी
        समीकरण $x'' + \omega_0^2 x = A\cos(\omega_0 t)$ को हल करता है
        (\cref{ex:b1:diffeq:oscillation} से तुलना कीजिए): \omterm{ex:b1:diffeq:oscillation}{अनुनाद} निरंतर धीमे
        और निरंतर बड़े विस्पंदों की सीमा है।
  \item \omterm{ex:b1:diffeq:oscillation}{अनुनाद} की दोनों नियतियों की तुलना कीजिए: अवमंदन के बिना रैखिक वृद्धि
  $\frac{At}{2\omega_0}$, और दुर्बल अवमंदन के साथ $R_{\max} \approx
  Q\,\frac{A}{\omega_0^2}$ पर संतृप्ति। एक वाक्य में: कौन-सी भौतिक
  क्रियाविधि पहली को दूसरी में बदल देती है?
\end{enumerate}

\medskip \noindent\textbf{भाग V --- ऊर्जा-संतुलन और संश्लेषण।}
\begin{enumerate}[resume]
  \item भाग II की स्थायी अवस्था में एक आवर्तकाल $\frac{2\pi}\Omega$ पर (क)
  प्रणोदन द्वारा प्रविष्ट शक्ति $P_{\mathrm{in}}(t) = A\cos(\Omega t) \cdot
  x_p'(t)$ और (ख) अवमंदन द्वारा क्षयित शक्ति $P_{\mathrm{diss}}(t) =
  2\lambda\,x_p'(t)^2$ का औसत संगणित कीजिए। दिखाइए कि दोनों औसत
  $\lambda\,R^2\Omega^2$ के बराबर हैं: प्रणोदन ठीक उतना ही देता है जितना
  अवमंदन जलाता है --- इसीलिए स्थायी अवस्था स्थायी है।
  \item इस समस्या में ठीक कहाँ प्रयोग हुआ: (क) संरचना प्रमेय
  \cref{thm:b1:diffeq:structure2}; (ख) सम्मिश्र चरघातांकी विधि; (ग)
  \cref{ch:b1:functions} की शैली में वास्तविक चर वाले फलन का अध्ययन?
  प्रत्येक के लिए एक वाक्य।
  \item संश्लेषण: $(E_\Omega)$ का पूरा व्यवहार-मानचित्र वर्णित कीजिए ---
  मुक्त बनाम प्रणोदित, अवमंदित बनाम अनवमंदित, और शिखर की ऊँचाई, पट्ट-चौड़ाई
  तथा क्षणिक भाग के जीवनकाल को एक साथ नियंत्रित करने वाले एकमात्र विमाहीन
  घुंडी $Q$ की भूमिका --- और बताइए कि कहानी आगे कहाँ चलती है: प्रथम कोटि के
  $2 \times 2$ निकाय (\cref{ch:b1:matrices} और स्नातक वर्ष 2 का खंड) तथा
  किसी व्यापक आवर्ती प्रणोदन का ज्यावक्रीय फलनों में वियोजन (फूरिये श्रेणी,
  स्नातक वर्ष 3 के खंड में), जिसके लिए इस समस्या की ज्यावक्रीय स्थिति मूल
  निर्माण-खंड है।
\end{enumerate}
\end{problem}
